% Tapis de Quillen, with Quillen's offprint and preprint (cote 162-5, typed part): archive cote 162-5, site page 50 (PDF page 51).
% offprint. Draft compared: none.
% AI-transcribed by Claude Opus 5.5 (claude-opus-5-5, adaptive thinking, effort high), Message Batches API, batch t162_full (msgbatch_01UAh3zyWPd454AeDkhnmQ33), 2026-11-13; not yet checked by eye.
\scanpage{162-5/50}
Inventiones math. 5, 237 -- 242 (1968)

On the Representation of Hermitian Forms
as Sums of Squares

Daniel G. Quillen* (Cambridge, Mass.)

There is an extensive literature (e.g. [1], Ch.VI) concerned with
Hilbert's problem of representing a positive-valued form with real
coefficients as a sum of squares. The corresponding problem for her-%
mitian forms does not seem to have been touched. In this paper we
present a theorem about hermitian forms which arose in connection
with a priori estimates for partial differential equations. This theorem is
proved by the methods of a priori estimates and one of its corollaries
is the Hilbert Nullstellensatz over $\mathbf{C}$, so we obtain a direct analytic proof
of the latter theorem.

Let $T$ and $V$ be finite dimensional complex vector spaces and let
$E(u_1,u_2)$ be a hermitian form on $S_m T\otimes V$ where $S_m$ denotes $m$-th sym-%
metric tensors. If $\zeta\in T$ and $v\in V$ let
\[ E^{\wedge}(\zeta,v)=(m!)^{-2}\,E(\zeta^m\otimes v,\zeta^m\otimes v). \]
Then $E^{\wedge}$ is a real-valued smooth function on $T\times V$ such that
\[ E^{\wedge}(s\,\zeta,t\,v)=|s|^{2m}\,|t|^2\,E^{\wedge}(\zeta,v) \qquad (1) \]
for $s,t\in\mathbf{C}$, and the map $E\mapsto E^{\wedge}$ yields a 1-1 correspondence between
hermitian forms on $S_m T\otimes V$and smooth functions on $T\times V$satisfying (1).
The inverse procedure of passing from a function to a form is polariza-%
tion (see formulas (2) below).

Let $|\zeta|^2$ be a hermitian scalar product on $T$ and define the ``$N$-th
prolongation'' of $E$ to be the hermitian form $E^N$ on $S_{m+N}T\otimes V$ with
\[ (E^N)^{\wedge}(\zeta,v)=(N!)^{-1}\,|\zeta|^{2N}\,E^{\wedge}(\zeta,v). \]
Theorem. The following conditions are equivalent:
(i) $E^{\wedge}(\zeta,v)>0$ for $\zeta\neq 0$ and $v\neq 0$.
(ii) For sufficiently large $N$, $E^N$ is positive definite.

We first derive an expression for $E^N$ as an integral. By choosing an
orthonormal basis for $T$ we may identify $T$ with the space of $\mathbf{C}$-linear
functions on $\mathbf{C}^n$ in such a way that the orthonormal basis is $z^1,\dots,z^n$
where $z^i$ is the $i$-th coordinate function. Then $S_k T$ may be naturally
identified with the holomorphic polynomial functions homogeneous of

* This research was supported in part by NSF GP-6959.

17 Inventiones math., Vol. 5

\mspage{\hand{50}}
